\scanpage{38}
On procède de même pour $f$, en notant \inserted{comme dessus} qu'il ne peut exister \uline{qu'un seul $f$} : \deleted{\unsure{défini}} \inserted{$P^1_S\to P^1_S$} définissant un revêtement d'ordre donné (ici 3) et dont la valeur en une section (p.ex. 0) disjointe de la ramification --- \unsure{normalisée en} deux sections disjointes --- soit donnée.

\unsure{En fait}, on ramène cette ramification à être $0,\infty$ et arrive au résultat
\[
s=f(t)=\frac{1+\left(\frac{\deleted{\unsure{\lambda}}t-x'}{\deleted{\unsure{\lambda}}t-x''}\right)^3}{1-\left(\frac{\deleted{\unsure{\lambda}}t-x'}{\deleted{\unsure{\lambda}}t-x''}\right)^3}
=\frac{(t-x'')^3+(t-x')^3}{(t-x'')^3-(t-x')^3}
\]
\margin{
\diagram{\begin{tikzcd}[column sep=huge,row sep=large]
\substack{\infty\\(x',x'')\deleted{\unsure{x}}\\P^1} \arrow[r,leftrightarrow,"{u^*(t)=\frac{t-x'}{t-x''}}","u"'] \arrow[d,"f"'] & \substack{1\\(0,\infty)\\P^1} \arrow[d,"{h^*(t)=-t^3}"] \\
\substack{P^1\\(1,-1)\\\infty} \arrow[r,"v","{v^*(t)=\frac{1-t}{1+t}}"'] & \substack{P^1\\(0,\infty)\\-1} \\
{} & {} \arrow[l,"{v^{-1*}(t)=\frac{1-t}{1+t}}"]
\end{tikzcd}}
\[
\begin{aligned}
f&=v^{-1}hu\\
f^*(t)&=u^*h^*v^{-1*}(t)=u^*h^*\left(\frac{1-t}{1+t}\right)\\
&=u^*\frac{1-h^*t}{1+h^*t}=u^*\frac{1+t^3}{1-t^3}\\
&=\frac{1+u^*(t)^3}{1-u^*(t)^3}=
\end{aligned}
\]
}
\figure{\begin{tikzpicture}[x=.7cm,y=.7cm]\draw (0,5)--(3.2,4.85) .. controls (3.5,4.9) and (3.25,4.6) .. (3.7,4.45) .. controls (4,4.25) and (3.6,4.15) .. (4.05,4.1) .. controls (4.3,4) and (4.9,3.5) .. (4.8,3.2) .. controls (4.7,2.4) and (5.05,1.8) .. (5,1.2)--(5,0);\end{tikzpicture}}
\[
=\frac{2t^3-3t^2(x''^{\deleted{\unsure{2}}}+x'^{\deleted{\unsure{2}}})+3t^{\deleted{\unsure{2}}}(x'^2+x''^2)-(x'^3+x''^3)}{3t^2(x'-x'')+3t(x''^2-x'^2)-(x''^3-x'^3)}
\]
Or ($x'$ et $x''$ étant racines \uline{primitives} 6ièmes de 1)
\[
x'^3=x''^3=-1,\quad\text{d'où }x''^3-x'^3=0
\]
\[
\left\{\begin{array}{l}x'+x''=1\\x'x''=1\end{array}\right.
\]
\[
\begin{aligned}
x'^2+x''^2&=(x'+x'')^2-2x'x''=-1\\
(x'-x'')^2&=(x'+x'')^2-4x'x''=-3\\
x'-x''&=\sqrt{-3}\\
x''^2-x'^2&=(x''-x')(x''+x')=-\sqrt{-3}\\
\deleted{\unsure{x''^3-x'^3}}&\deleted{\unsure{=(x''-x')(x''^2+x'x''+x'^2)=0}}
\end{aligned}
\]
\[
\begin{aligned}
f(t)&=\frac1{3\sqrt{-3}}\frac{2t^3-3t^2-3t+2}{t(t-1)}\\
&=\frac1{3\sqrt{-3}}\frac{(t+1)(2t-1)(t-2)}{t(t-1)}
\end{aligned}
\]
\mspage{37}
